\chapter{Introduction}\label{r03:ch:intro}

\emph{What problem this report addresses, the goal and the rules of comparison, how the work
was done, and the result in brief. Every technical term used here is explained in
\cref{r03:ch:background} and listed in the glossary (Appendix~\ref{r03:app:glossary}).}

\section{The problem}\label{r03:sec:problem}

waLBerla is a C++ framework for fluid simulations with the lattice Boltzmann method (LBM).
Its compute kernels are not written by hand: two Python packages, lbmpy and \ps{}, derive the
update formulas and generate C++ code from them (\cref{r03:sec:bg-stack}).

The SpacemiT~K1 is an eight-core RISC-V system on a chip, used here on a Banana~Pi BPI-F3
board. Its cores, the SpacemiT~X60, have a vector unit: the RISC-V Vector extension (RVV),
which processes 8 single-precision or 4 double-precision numbers with one instruction
(\cref{r03:sec:bg-rvv}). Earlier work added an RVV code generator to \ps{} and measured
waLBerla's three-dimensional LBM kernel (D3Q19) with it on one X60 core.\footnote{Report~2 of
this collection. Its numbers that this report starts from are restated here, so that report
is not needed to follow this one.} The best vector version of the program was only
\sx{1.73} faster than the same program without vector instructions in single precision
(GCC; \sx{1.92} with Clang), and \sx{1.31} faster in double precision (GCC), on a grid of
$64^3$ cells. At that time the gap to the possible \sx{8} and \sx{4} was attributed to
memory bandwidth.

This report follows the work that came after. It answers three questions:
\begin{enumerate}
  \item What really limits the vectorized kernel on the X60?
  \item Which changes to the generated code, to the code generators and to waLBerla remove
        that limit, and by how much?
  \item How close does the complete application get to \sx{8} (single precision) and
        \sx{4} (double precision)?
\end{enumerate}

\section{Goal and rules of comparison}\label{r03:sec:goal}

\paragraph{Goal.} The goal was a speedup of \sx{8} in single precision (f32) and \sx{4} in
double precision (f64). The speedup is measured in the complete benchmark application, on
one core, for a grid of $64^3$ cells. It is the ratio of the vector version's throughput to
that of the \term{plain scalar build}: the same application compiled without vector
instructions. Throughput is given in MLUP/s, million lattice-cell updates per second.
\Cref{r03:sec:ops} shows that \sx{8} and \sx{4} are the real arithmetic limits of the
X60: its vector unit does 8 (f32) or 4 (f64) operations in the time its scalar unit does one.

\paragraph{Fair baseline.} Some of the changes in this report are not specific to vector
code. Examples are common-subexpression elimination, in-place streaming, the data layout
and huge memory pages. A scalar program can use them as well. So that such changes do not
inflate the speedup, they were also applied to the scalar program wherever possible. This
gives a second baseline, the \term{best scalar build} (\cref{r03:tab:variants}). Every
result is reported in three ways: the absolute throughput in MLUP/s, the ratio to the plain
scalar build, and the ratio to the best scalar build.

\paragraph{One core.} The K1 is an embedded system on a chip, and all results are for one
core. A short side note on several cores is in \cref{r03:sec:res-mpi}.

\section{How the work was done}\label{r03:sec:method}

The work proceeded in five steps:
\begin{enumerate}
  \item \textbf{Measure} where the existing program spends its time (\cref{r03:ch:root-cause}).
  \item \textbf{Characterise} the X60 core with small test programs (\term{probes}) that
        measure memory and instruction costs (\cref{r03:sec:mem,r03:sec:ops}).
  \item \textbf{Try} each candidate optimisation in a \term{kernel lab}: a small program
        that times one kernel outside waLBerla, so that a change can be tested in minutes
        (\cref{r03:sec:lab,r03:ch:levers}).
  \item \textbf{Move} the optimisations that pay off into \ps{}, sweepgen and waLBerla
        (\cref{r03:ch:pystencils,r03:ch:walberla}), and check that the results stay
        correct (\cref{r03:ch:verification}).
  \item \textbf{Measure} the complete application again (\cref{r03:ch:results}).
\end{enumerate}

Every optimisation ends in one of three states, shown in the same colours in all figures:
\begin{description}
  \item[Adopted] (teal) in the committed code of \cref{r03:tab:commits}, and measured;
  \item[Rejected] (red) measured, no gain, not in the code;
  \item[Pending] (orange) in the code or in the lab, but the measurement that decides it was
        still running when this report was written (\cref{r03:ch:open}).
\end{description}
Findings about the hardware, which are not changes, are shown in blue.

\section{Result in brief}\label{r03:sec:brief}

\begin{keyidea}[Where the application stands]
waLBerla's benchmark application \code{SweepgenLbmCPU}, $64^3$ cells, one core, GCC~14, the
whole program (kernel plus data exchange), best of 3 runs of 200 time steps. Latest complete
measurement (run~C of \cref{r03:sec:res-runs}):
\begin{itemize}
  \item \textbf{single precision: 7.25~MLUP/s, \sx{5.24} the plain scalar build}
        (1.38~MLUP/s), and \sx{3.54} the best scalar build (2.05~MLUP/s);
  \item \textbf{double precision: 4.04~MLUP/s, \sx{4.23} the plain scalar build}
        (0.95~MLUP/s), and \sx{2.78} the best scalar build (1.45~MLUP/s).
\end{itemize}
The \sx{4} goal for double precision is met, with generated code. The \sx{8} goal for single
precision is not met yet (\cref{r03:sec:gap}).
\end{keyidea}

The limit was not memory bandwidth (\cref{r03:ch:root-cause}). The X60 has a hardware
prefetcher: a unit that detects regular access patterns and loads data before the program
asks for it. It can follow about 8 such patterns (\term{streams}) at a time. The kernel
had 42 (19 read and 23 write streams), so almost every access that missed the cache
waited for the full memory latency. Three optimisations attacked this and were adopted:
\begin{enumerate}
  \item \textbf{software prefetching}, a new option of \ps{}: the generated loop requests
        its data 128~bytes ahead (\cref{r03:sec:ps-prefetch}); about \sx{2} on its own;
  \item \textbf{in-place streaming (AA pattern)}, in sweepgen and in waLBerla's data
        exchange: one data array instead of two, 23 streams instead of 42
        (\cref{r03:sec:wl-sweepgen,r03:sec:wl-halo});
  \item \textbf{a better RVV instruction selection} in \ps{}: scalar operands and fused
        multiply-add, which removed all register spills from the loop
        (\cref{r03:sec:ps-selector}).
\end{enumerate}
A fourth adopted change computes each row's addresses once and keeps constants in
registers (\cref{r03:sec:ps-rowptr}). It adds 2.5\,\% (f32) and 2.2\,\% (f64).

Three more changes are in the code, but their end-to-end measurement was still running
(\cref{r03:ch:open}):
\begin{itemize}
  \item a faster data exchange between neighbouring blocks, which moves one data slice at a
        time (\cref{r03:sec:wl-halo});
  \item time steps that skip writing the density and velocity fields
        (\cref{r03:sec:wl-sweepgen});
  \item a kernel that splits the update into passes over a few directions at a time,
        generated from lbmpy (\cref{r03:sec:wl-split}).
\end{itemize}

\begin{figure}[htbp]
  \centering
  \resizebox{\linewidth}{!}{% Map of the work: the findings about the hardware on top (in the order they were made), the
% optimisations below, grouped by what they attack and coloured by outcome. Numbers in brackets
% are sections of this report.
\begin{tikzpicture}[
    every node/.style={font=\scriptsize},
    lev/.style={text width=34mm,minimum height=10mm,font=\scriptsize},
  ]
  % ---- findings -------------------------------------------------------------------
  \node[r03info,lev] (f0) at (0,0)     {\textbf{not bandwidth}: 0.4~GB/s, flat over grid sizes (\ref{r03:sec:bw})};
  \node[r03info,lev] (f1) at (3.75,0)  {\textbf{latency bound}: 2.97 DRAM lines per cell, prefetcher idle (\ref{r03:sec:cycles})};
  \node[r03info,lev] (f2) at (7.5,0)   {\textbf{too many streams}: prefetcher follows $\approx$\,8, the kernel has 42 (\ref{r03:sec:mem})};
  \node[r03info,lev] (f5) at (11.25,0) {\textbf{limit} \sx{8} f32, \sx{4} f64: 1 vector instruction per cycle (\ref{r03:sec:ops})};
  \draw[r03arrow] (f0) -- (f1); \draw[r03arrow] (f1) -- (f2); \draw[r03arrow] (f2) -- (f5);

  % ---- the four kinds of optimisation ---------------------------------------------------
  \foreach \x/\t [count=\i] in {0/start loads earlier, 3.75/fewer memory streams,
                                7.5/shorter vector loop, 11.25/time outside the kernel} {
    \node[font=\footnotesize\bfseries,text=psblue!85!black] (h\i) at (\x,-1.3) {\t};
  }
  \draw[psgray!60] (-1.75,-0.95) -- (13.0,-0.95);
  \node[r03label,anchor=east] at (13.0,-0.78) {optimisations, by what they attack};

  % column 1: prefetch
  \node[r03ok,lev] (pf) at (0,-2.45) {\textbf{software prefetch}\\128~B ahead (\ref{r03:sec:ps-prefetch})\\lab \sx{2.1}, application \sx{1.96}};
  % column 2: streams
  \node[r03box,lev] (sp) at (3.75,-2.45) {direction split, hand kernel: 6.98 vs 5.2~MLUP/s (\ref{r03:sec:lever-split})};
  \node[r03box,lev] (aah) at (3.75,-4.0) {AA + split, hand kernel: 180 cycles per cell (\ref{r03:sec:lever-aa})};
  \node[r03ok,lev] (aag) at (3.75,-5.55) {\textbf{AA streaming} from lbmpy, in sweepgen and the halo exchange (\ref{r03:sec:wl-sweepgen}, \ref{r03:sec:wl-halo})};
  \node[r03wait,lev] (spg) at (3.75,-7.1) {direction split, generated from lbmpy (\ref{r03:sec:wl-split})};
  \draw[r03arrow] (sp) -- (aah); \draw[r03arrow] (aah) -- (aag); \draw[r03arrow] (aag) -- (spg);
  % column 3: the loop body
  \node[r03no,lev] (lm) at (7.5,-2.45) {LMUL 2 / 4: slower, 216 vs 163 cycles per cell (\ref{r03:sec:lever-sched})};
  \node[r03ok,lev] (tu) at (7.5,-4.0) {X60 scheduling model, Clang only (\ref{r03:sec:wl-cmake})};
  \node[r03ok,lev] (vf) at (7.5,-5.55) {\textbf{.vf operands + FMA}\\lab \sx{1.22} (\ref{r03:sec:ps-selector})};
  \node[r03ok,lev] (rp) at (7.5,-7.1) {\textbf{row pointers}, constants in registers: +2.5\,\% f32, +2.2\,\% f64 (\ref{r03:sec:ps-rowptr})};
  % column 4: around the kernel
  \node[r03no,lev] (al) at (11.25,-2.45) {alignment 32/64~B, huge pages: 16~B is best (\ref{r03:sec:lever-layout})};
  \node[r03wait,lev] (ha) at (11.25,-4.0) {halo exchange one slot at a time; 21\,\% of a step (\ref{r03:sec:wl-halo})};
  \node[r03wait,lev] (pa) at (11.25,-5.55) {padding between the PDF slices (\ref{r03:sec:lever-layout})};
  \node[r03wait,lev] (no) at (11.25,-7.1) {no density/velocity output in the timed steps (\ref{r03:sec:wl-sweepgen})};

  % ---- result ----------------------------------------------------------------------------
  \node[r03ok,text width=140mm,minimum height=9mm] (res) at (5.625,-8.65)
    {\footnotesize\textbf{Application, latest run (run C):} f32 7.25~MLUP/s $=$ \sx{5.24}, f64 4.04~MLUP/s $=$
     \sx{4.23} the plain scalar build (\sx{3.54} / \sx{2.78} the best scalar build)};
\end{tikzpicture}
}
  \caption[Map of the findings and optimisations]{Map of the work. Top: the findings about
    the hardware, in the order they were made. Below: the optimisations, grouped by what they
    attack, coloured by outcome (teal adopted, red rejected, orange pending; white: lab
    experiments that led to an adopted change). Numbers in brackets are sections of this
    report.}\label{r03:fig:story}
\end{figure}

\section{Where the code is}\label{r03:sec:code}

All work is committed and pushed to GitHub. Every link in this report points at these
commits:

\begin{table}[htbp]
  \centering\small
  \caption{Repositories and commits this report describes.}\label{r03:tab:commits}
  \begin{tabularx}{\linewidth}{@{}lllX@{}}
    \toprule
    \textbf{Repository} & \textbf{Branch} & \textbf{Commit} & \textbf{Contents}\\
    \midrule
    \href{\rWSurl}{\texttt{my-e4-ws}} & \texttt{main} & \rcommitws &
      the workspace: working notes and raw measurement logs (\file{docs/notes3/}), kernel lab
      and probes (\file{misc/x60-lab/}), benchmark scripts (\file{scripts/})\\
    \href{\rPSurl}{\texttt{mirrored-pystencil2}} & \texttt{rvv-support} & \rcommitps &
      \ps{} with the RVV code generator: prefetch option, instruction selection, row
      pointers. Changes described here:
      \href{\rPSurl/compare/9a989e1625fbb0272876f550a5f7d3456fdaca25...\rPSsha}{\texttt{9a989e16..a81dc547}}\\
    \href{\rWLurl}{\texttt{mirrored-walberla}} & \texttt{rvv-support} & \rcommitwl &
      waLBerla with sweepgen: in-place streaming, data exchange, benchmark application,
      build flags. Changes described here:
      \href{\rWLurl/compare/b78f2d8791589fa38d0e407a3447865cb8ce227f...\rWLsha}{\texttt{b78f2d87..efa2b70c}}\\
    \bottomrule
  \end{tabularx}
\end{table}

The workspace commit \rcommitws{} includes both forks as Git submodules, at exactly these
commits. The copies of the raw measurement logs (Appendix~\ref{r03:app:evidence}) and some
documentation edits are newer than \rcommitws: they are in the working tree, not yet in a
commit.

\section{How to read this report}\label{r03:sec:guide}

\begin{description}
  \item[Part I, Background and Diagnosis] \Cref{r03:ch:background} explains the method, the
        software and the hardware terms. \Cref{r03:ch:setup} describes the board, the
        benchmark application and the kernel lab. \Cref{r03:ch:root-cause} shows why the
        vectorized kernel stalls.
  \item[Part II, Optimisations] \Cref{r03:ch:levers} tests every candidate in the kernel
        lab. \Cref{r03:ch:pystencils} explains how the generated code is vectorized and what
        changed in \ps{}; \cref{r03:ch:walberla} what changed in sweepgen and waLBerla.
  \item[Part III, Evidence] \Cref{r03:ch:verification} shows how each change was checked
        for correctness, \cref{r03:ch:results} what the application reaches, and
        \cref{r03:ch:open} what was still being measured.
\end{description}
Appendix~\ref{r03:app:glossary} lists every term and abbreviation,
appendix~\ref{r03:app:file-map} every file that changed, appendix~\ref{r03:app:repro} how
to reproduce the results, and appendix~\ref{r03:app:evidence} the log file behind each
number.

\paragraph{Notation.} f32 and f64 stand for single precision (32-bit floating point) and
double precision (64-bit). A factor such as \sx{5.24} is a speedup: the throughput of one
version divided by that of another. ``Cycles per cell'' is the number of processor clock
cycles spent per cell update; at the X60's 1.6~GHz, cycles per cell $= 1600 /
\text{MLUP/s}$.
